% ============================================================
% Appendice B — Comandi conversazionali completi
% ============================================================

\chapter{Comandi Conversazionali Completi}
\label{app:comandi}

Riferimento completo di tutti i comandi conversazionali AI-DLC.

\section*{Tabella di riferimento rapido}

\begin{table}[htbp]
  \centering
  \begin{tabular}{lll}
    \toprule
    \textbf{Comando} & \textbf{Scopo} & \textbf{Bypass sicurezza?} \\
    \midrule
    \texttt{@Task [obiettivo]} & Crea e avvia un task (test-driven) & No \\
    \texttt{@checkpoint} & Fissa stato avanzamento & No \\
    \texttt{@context-update} & Propone aggiornamento context & No \\
    \texttt{@show-constraints} & Mostra vincoli SEC/PERF & No \\
    \texttt{@security-check} & Verifica codice vs SEC & No \\
    \texttt{@perf-check} & Verifica codice vs PERF & No \\
    \texttt{@load-phase [N]} & Carica modulo fase N & No \\
    \texttt{@set-phase [N]} & Propone cambio fase & No \\
    \texttt{@prompt-list} & Lista prompt riusabili & No \\
    \texttt{@prompt [ID]} & Applica template prompt & No \\
    \texttt{@alternatives} & Genera 2--3 opzioni & No \\
    \texttt{@simplify} & Riduce scope/complessità & No \\
    \texttt{@rollback} & Piano di rollback & No \\
    \texttt{@stop} & Ferma immediatamente & No \\
    \bottomrule
  \end{tabular}
  \caption{Riferimento rapido comandi AI-DLC}
  \label{tab:app-b-comandi}
\end{table}

\section*{\texttt{@Task} --- avvio dell'implementazione}

Crea l'unità di implementazione e ne avvia il ciclo test-driven (Capitolo~\ref{ch:fasi-e-modes}, §\ref{sec:06-task-tdd}). Input: un obiettivo o un riferimento a un Epic/requisito. Produce il task \texttt{T-NNN} con AI Sizing, criteri di accettazione e i test che li codificano (da scrivere prima del codice). Per i task MEDIUM o superiori il piano resta in attesa di approvazione; dopo l'approvazione l'agente implementa finché i test del task passano, poi esegue l'intera suite (test del task + regressione) e propone il \texttt{@checkpoint}. Rispetta sempre HALT trigger e risk floor.

\section*{Regole di tutti i comandi}

\begin{enumerate}
  \item Nessun comando bypassa gli HALT trigger. Neanche \texttt{@stop} li disattiva.
  \item Nessun comando modifica file automaticamente --- propone sempre, tu applichi.
  \item \texttt{@security-check} e \texttt{@perf-check} citano solo i vincoli in \texttt{\_CONTEXT.md}.
  \item \texttt{@rollback} non esegue azioni distruttive senza conferma esplicita.
  \item I comandi non cambiano la classificazione del rischio.
\end{enumerate}

\section*{Comandi personalizzati di progetto}

In \texttt{.ai-dlc/project/instructions.md}, definisci shortcut per workflow ricorrenti:

\begin{lstlisting}[language=,caption={Formato comandi personalizzati}]
## Project Commands

@nome-comando
  [Descrizione di cosa deve fare l'agente quando
   riceve questo comando]
\end{lstlisting}

Vedi Capitolo~\ref{ch:comandi-conversazionali} per esempi e output attesi di ciascun comando.
